\subsection{子环}

定义 3. 设 A 是一个环。A 的子环是 A 的任意子集 B，它在加法下是 A 的子群，在乘法下稳定，并且包含 A 的单位元。

上述条件可以写成如下形式

\[
0 \in \mathbf {B}, \quad \mathbf {B} + \mathbf {B} \subset \mathbf {B}, \quad - \mathbf {B} \subset \mathbf {B}, \quad \mathbf {B}. \mathbf {B} \subset \mathbf {B}, \quad 1 \in \mathbf {B}.
\]

若 B 是 A 的子环，则给它赋予由 A 上的加法和乘法诱导的运算，这使它成为一个环。从 B 到 A 的典范单射是一个环同态。

例子。(1) 加法群 \(\mathbf{Z}\) 的每个包含 1 的子群都等于 \(\mathbf{Z}\) 。因此 \(\mathbf{Z}\) 是 \(\mathbf{Z}\) 的唯一子环。

\begin{enumerate}
\def\labelenumi{(\arabic{enumi})}
\setcounter{enumi}{1}
\item
  设 A 为一个环，且 \((A_{t})_{t\in I}\) 为 A 的一族子环；显然 \(\bigcap_{i\in I}A_{i}\) 是 A 的子环。特别地，A 的包含 A 的子集 X 的所有子环的交集是一个子环，称为由 X 生成的 A 的子环。
\item
  设 X 是环 A 的一个子集。X 在 A 中的中心化子是 A 的一个子环。特别地，A 的中心是 A 的一个子环。
\item
  设 G 是一个带算子的交换群；设 \(\Omega\) 表示算子集， \(\alpha \mapsto f_{\alpha}\) 表示 \(\Omega\) 在 G 上的作用。设 E 是不带算子的群 G 的自同态环，F 是带算子的群 G 的自同态组成的集合。根据定义，F 由 G 的自同态 \(\phi\) （满足 \(\phi \cdot f_{\alpha} = f_{\alpha} \cdot \phi\) ，对所有 \(\alpha \in \Omega\) ）组成。因此 F 是环 E 的一个子环。F 称为带算子的群 G 的自同态环（参见 II，§ 1，第 2 号）。设 \(F_{1}\) 是由 \(f_{\alpha}\) 生成的 E 的子环。则 F 是 \(F_{1}\) 在 E 中的中心化子。
\end{enumerate}

